\section{Bus Interface Unit: Bus, Queue, dan Alamat}
\label{sec:biu}

\eng{Bus Interface Unit} (BIU) adalah pintu keluar-masuk prosesor.
Seluruh instruksi dan data yang berasal dari memori atau port harus
melewati unit ini. Bagian ini diringkas dari \cite{gunawan2019},
Bagian 2, hlm.~12--15, dan \cite{sto2001}, BAB II, hlm.~9--10.

\subsection{Tugas Bus Interface Unit}

Menurut buku sumber, BIU mempunyai tugas utama sebagai berikut.

\begin{enumerate}
    \item mengirim tugas ke bagian lain di dalam prosesor;
    \item mengambil instruksi dari memori;
    \item membaca atau menulis data dari dan ke memori serta port;
    \item menangani pengalamatan dan transfer data.
\end{enumerate}

Dengan kata lain, BIU mengurus seluruh hubungan antara prosesor dan
dunia luar. EU yang menjalankan instruksi tidak perlu mengetahui
bagaimana sinyal bus bekerja; EU cukup meminta data atau alamat
kepada BIU.

\subsection{Komponen Bus Interface Unit}

BIU memiliki empat komponen utama yang saling berkaitan.

\begin{description}
    \item[Bus] Jalur tempat instruksi dan data mengalir antar-
          komponen. Bus menghubungkan prosesor dengan memori dan
          perangkat lain.

    \item[\eng{Queue}] Kumpulan register berukuran enam byte yang
          menjadi antrean instruksi sebelum masuk ke EU. Antrean ini
          diisi oleh BIU dan diambil oleh EU.

    \item[Register segmen] Empat register 16 bit, yaitu \reg{CS},
          \reg{DS}, \reg{SS}, dan \reg{ES}. Register ini menyimpan
          nomor segmen untuk kode, data, \eng{stack}, dan segmen
          tambahan.

    \item[\eng{Instruction Pointer} (IP)] Register 16 bit yang
          menyimpan \eng{offset} instruksi berikutnya. Pasangan
          \reg{CS}:\reg{IP} menunjukkan lokasi instruksi yang akan
          diambil.
\end{description}

\subsection{Instruction Queue Enam Byte}

Memori bekerja lebih lambat daripada prosesor. Agar prosesor tidak
selalu menunggu instruksi datang, BIU mengisi \eng{queue} dengan
instruksi berikutnya sambil EU mengerjakan instruksi yang sekarang.
\eng{Queue} pada 8086 berukuran enam byte dan akan diisi kembali
ketika isinya berkurang dua byte karena diteruskan ke ALU.

Tanpa antrean, pengambilan instruksi dan eksekusi berjalan
bergantian. Dengan antrean, keduanya dapat berjalan bersamaan
sebagaimana ditunjukkan pada Tabel~\ref{tab:queue}.

\begin{table}[H]
    \centering
    \caption{Perbandingan pengambilan dan eksekusi tiga instruksi
    tanpa dan dengan \eng{instruction queue}.}
    \label{tab:queue}
    \small
    \begin{tabular}{@{}cll@{}}
        \toprule
        \textbf{Siklus} & \textbf{Tanpa Antrean} &
        \textbf{Dengan Antrean} \\
        \midrule
        1 & ambil instruksi 1 & ambil instruksi 1 \\
        2 & eksekusi instruksi 1 & eksekusi 1 dan ambil 2 \\
        3 & ambil instruksi 2 & eksekusi 2 dan ambil 3 \\
        4 & eksekusi instruksi 2 & eksekusi 3 \\
        5 & ambil instruksi 3 & --- \\
        6 & eksekusi instruksi 3 & --- \\
        \bottomrule
    \end{tabular}
\end{table}

Tabel~\ref{tab:queue} adalah model sederhana, bukan angka waktu
sebenarnya. Intinya, antrean memungkinkan pengambilan instruksi
berikutnya berjalan sementara EU mengeksekusi instruksi sekarang.

\begin{peringatanbox}
\vspace{1.1em}
\eng{Queue} hanya berukuran enam byte. Bila instruksi yang sedang
dijalankan membutuhkan operand dari memori, BIU harus menghentikan
sementara pengisian antrean dan melayani permintaan operand tersebut.
Karena itu, program yang banyak mengakses memori tidak memperoleh
keuntungan antrean sebesar program yang bekerja di register.
\end{peringatanbox}

\subsection{Tiga Kelompok Bus}

BIU berhubungan dengan memori dan port melalui tiga kelompok bus.
Pembagiannya dirangkum pada Tabel~\ref{tab:bus}.

\begin{table}[H]
    \centering
    \caption{Tiga kelompok bus pada mikroprosesor 8086.}
    \label{tab:bus}
    \small
    \begin{tabular}{@{}p{2.7cm}p{1.8cm}p{2.4cm}p{6.0cm}@{}}
        \toprule
        \textbf{Bus} & \textbf{Lebar} & \textbf{Arah} &
        \textbf{Fungsi} \\
        \midrule
        Bus alamat & 20 bit & keluar & Membawa alamat lokasi memori
        atau port yang akan diakses \\
        \addlinespace
        Bus data & 16 bit & dua arah & Membawa data yang dibaca atau
        ditulis \\
        \addlinespace
        Bus kontrol & sinyal & keluar/masuk & Mengatur jenis, arah,
        dan waktu pertukaran data \\
        \bottomrule
    \end{tabular}
\end{table}

Bus alamat 8086 selebar 20 bit sehingga ruang alamat yang dapat
dijangkau adalah $2^{20}$ byte atau 1 MB. Bus data selebar 16 bit
berarti prosesor dapat memindahkan dua byte sekaligus. Bus kontrol
menentukan apakah operasi yang sedang berlangsung adalah membaca
atau menulis, serta kapan perangkat luar siap melayani.

\subsection{Perhitungan Alamat Fisik}

Register segmen hanya 16 bit, sedangkan bus alamat 20 bit. BIU
menjembatani keduanya dengan rumus:

\begin{equation}
    \text{alamat fisik} = (\text{nilai segmen} \times 16)
    + \text{offset}.
    \label{eq:alamat-fisik}
\end{equation}

Perkalian dengan 16 sama dengan menggeser nilai segmen satu digit
heksadesimal ke kiri. Contoh, bila \reg{CS} berisi \hex{1234} dan
\reg{IP} berisi \hex{0010}, maka instruksi berikutnya diambil dari
alamat fisik:

\[
    \hex{1234} \times 16 + \hex{0010}
    = \hex{12340} + \hex{0010} = \hex{12350}.
\]

Perhitungan ini dikerjakan oleh BIU. Pembahasan lebih rinci mengenai
segmentasi dan \eng{overlapping} disajikan pada bab tentang
segmentasi memori.

\begin{infobox}
\vspace{1.1em}
Pasangan \reg{CS}:\reg{IP} selalu menunjuk instruksi berikutnya.
Setiap kali BIU mengambil satu byte instruksi, \reg{IP} bertambah
sesuai panjang byte tersebut. Perubahan \reg{IP} inilah yang dapat
diamati ketika program dijalankan dengan \eng{single step}.
\end{infobox}
